% !TeX root = ../thesis.tex
\chapter{LÝ THUYẾT}


Trong quá trình nghiên cứu và thực hiện đề tài, chúng em đã tham khảo nhiều công trình nghiên cứu và bài báo khoa học trên thế giới \cite{siciliano2009, craig2005, gonzalez2018}. Tuy nhiên, tại Việt Nam, các tài liệu và báo cáo khoa học đi sâu vào việc hiện thực hóa toàn diện hệ thống điều khiển cho robot bốn chân vẫn còn khá hạn chế. Do đó, chương này sẽ tổng hợp các cơ sở lý thuyết nền tảng cốt lõi nhất làm tiền đề cho việc thiết kế và điều khiển hệ thống.

\section{Tổng quan về Robot bốn chân}
Robot bốn chân là một hệ thống robot di động sử dụng cơ cấu bốn chân, được thiết kế dựa trên nguồn cảm hứng sinh học từ đặc điểm vận động của các loài côn trùng và động vật có vú. Khác với robot hai chân, robot bốn chân có trọng tâm thấp và diện tích đa giác chân đế rộng hơn. Nhờ đó, chúng có thể dễ dàng tự ổn định thân mình bằng cách linh hoạt điều chỉnh tư thế cơ thể cũng như thay đổi các dáng đi. Với đặc tính cơ động học ưu việt này, robot bốn chân sở hữu khả năng thích nghi tuyệt vời trên các địa hình gồ ghề và phức tạp.

Trong khuôn khổ của nghiên cứu này, đối tượng nghiên cứu và chế tạo cụ thể là một nguyên mẫu robot bốn chân mô phỏng hình dáng và cấu trúc khớp của loài chó. Cấu trúc này mang lại sự cân bằng tối ưu giữa kích thước, sự linh hoạt và tải trọng, là nền tảng lý tưởng để triển khai các giải thuật điều khiển thăng bằng.

\begin{figure}[H]
    \centering
    \includegraphics[width=0.7\textwidth]{hardware/robot_overview.png}
    \caption{Ví dụ robot bốn chân}
    \label{fig1}
\end{figure}

\section{Phân loại cấu trúc chân robot}
Cấu trúc cơ khí của chân robot có ảnh hưởng lớn đến khả năng điều khiển, hiệu suất chuyển động, độ chính xác cũng như khả năng mở rộng trong thiết kế. Dựa trên nguyên lý cơ học, cấu trúc chân robot có thể được chia thành ba loại chính: cấu trúc nối tiếp (Serial), cấu trúc song song (Parallel) và cấu trúc lai (Hybrid).

\textbf{\textit{Cấu trúc nối tiếp (Serial)}} 

Là cấu trúc phổ biến nhất, trong đó các khớp được liên kết tuần tự từ gốc (khung thân robot) đến điểm cuối (bàn chân). Mỗi khớp điều khiển một bậc tự do (DOF) và ảnh hưởng trực tiếp đến toàn bộ các khâu phía sau trong chuỗi động học. Hình \ref{fig8} minh họa một dạng cấu trúc chân nối tiếp với ba động cơ được gắn lần lượt tại khớp vai, khớp hông và khớp gối:

\begin{figure}[H]
    \centering
    \includegraphics[width=0.6\textwidth]{hardware/serial_leg_both_sides.png}
    \caption{Cơ cấu chân dạng nối tiếp}
    \label{fig8}
\end{figure}

\textbf{\textit{Cấu trúc song song (Parallel)}} 

Trong cấu trúc này, nhiều nhánh cơ cấu hoạt động đồng thời để điều khiển một điểm cuối duy nhất, tạo nên hệ thống động học khép kín. Loại này thường có cấu hình tương tự như bệ Stewart (một dạng bệ lục giác có khung đồng bộ với các chân), mang lại độ cứng vững cao và độ chính xác chuyển động tốt hơn so với cấu trúc nối tiếp. Tuy nhiên, cấu trúc này có khối lượng cơ khí lớn và thuật toán tính toán động học thuận (FK), động học nghịch (IK) phức tạp hơn đáng kể. Hình \ref{fig9} trình bày một dạng cấu trúc chân song song phổ biến:

\begin{figure}[H]
    \centering
    \includegraphics[width=0.6\textwidth]{hardware/mechanical_model.png }
    \caption{Mô hình cơ khí và hệ tọa độ nhánh của chân song song}
    \label{fig9}
\end{figure}

\textbf{\textit{Cấu trúc lai (Serial-Parallel Hybrid)}} 

Là sự kết hợp ưu điểm của cả hai cấu trúc trên, trong đó hệ thống sử dụng chuỗi động học nối tiếp để điều khiển chuyển động chính, đồng thời tích hợp cơ cấu song song để hỗ trợ truyền động (thường ứng dụng để dẫn động từ xa cho khớp gối). Do kết hợp cả hai cơ chế, phần chân robot thường nặng nề về mặt cơ khí và thuật toán điều khiển chuyển động cũng trở nên vô cùng phức tạp. Hình \ref{fig10} mô phỏng một trong những cấu trúc chân lai đơn giản nhất:

\begin{figure}[H]
    \centering
    \includegraphics[width=0.6\textwidth]{hardware/mechanical_model_branch.png  }
    \caption{Cơ cấu chân lai nối tiếp – song song}
    \label{fig10}
\end{figure}

Sau quá trình đánh giá, nhằm tối ưu hóa trọng lượng phần cứng và tập trung vào phát triển giải thuật điều khiển cũng như thuận lợi cho việc kiểm chứng mô phỏng, cấu trúc chân nối tiếp (Serial) đã được lựa chọn cho nguyên mẫu robot.
\section{Cấu trúc liên kết động học của robot bốn chân}
Mặc dù tồn tại nhiều dạng cấu trúc liên kết động học khác nhau trong thiết kế robot bốn chân, các cấu hình với 8, 12 hoặc 16 bậc tự do (DOF) được sử dụng phổ biến nhất.

Cấu trúc 8 bậc tự do bao gồm bốn khớp hông và bốn khớp gối, trong đó tất cả các trục khớp song song với nhau. Robot 8 DOF có khả năng di chuyển tiến và lùi với tốc độ cao và tương đối dễ điều khiển do không có bậc tự do vung ngang tại khớp hông. Tuy nhiên, robot bốn chân 8 DOF không thể thực hiện chuyển động theo phương ngang và khả năng quay hướng bị hạn chế đáng kể, dẫn đến hiệu suất chuyển động không cao.

Ngược lại, robot bốn chân 16 DOF có số lượng khớp nhiều hơn, cho phép thực hiện các chuyển động linh hoạt và khéo léo hơn, nhưng đồng thời làm tăng đáng kể độ phức tạp trong điều khiển và tính toán.

Dựa trên những phân tích trên, cấu hình 12 bậc tự do, với mỗi chân gồm ba khớp quay, được lựa chọn cho thiết kế trong nghiên cứu này.

Về cấu hình chân, các dạng phổ biến bao gồm: inner knee-elbow, outer knee-elbow, full-knee và full-elbow. Hai cấu hình đầu có cấu trúc chân đối xứng hoàn toàn, do đó giúp robot đạt được độ ổn định cao hơn. Trong khi đó, hai cấu hình sau cung cấp không gian làm việc lớn hơn và thuận lợi hơn cho việc điều khiển chuyển động. Vì vậy, trong nghiên cứu này, cấu hình full-elbow được lựa chọn.

Cấu trúc liên kết động học kết cấu rút gọn được thể hiện ở Hình \ref{fig11}, trong đó khớp hông có hai bậc tự do quay (vung ngang và vung dọc), còn khớp gối có thêm một bậc tự do quay.

\begin{figure}[H]
    \centering
    \includegraphics[width=0.6\textwidth]{hardware/structure_topology_front.png }
    \caption{Cấu trúc liên kết động học của robot bốn chân 12 bậc tự do}
    \label{fig11}
\end{figure}

\section{Cơ cấu chân 3 bậc tự do}
Đối với robot bốn chân, cơ cấu chân 3 bậc tự do được ưu tiên sử dụng phổ biến nhất nhờ tính linh hoạt cao trong việc mô phỏng các dáng đi sinh học. Một chân robot tiêu chuẩn thường được chia thành ba khâu chính: hông (\textit{hip}), đùi (\textit{thigh}) và cẳng chân (\textit{calf} hoặc \textit{lower thigh}). 

\begin{figure}[H]
    \centering
    \includegraphics[width=0.6\textwidth]{kinematics/single_leg_3dof.png}
    \caption{Ba bậc tự do của cơ cấu một chân}
    \label{fig2}
\end{figure}

Như vậy, để cấu thành nên một robot bốn chân hoàn chỉnh, hệ thống sẽ cần tổng cộng 12 bậc tự do. Mỗi khớp tương ứng với một bậc tự do và được dẫn động độc lập bởi một động cơ. Sự phối hợp đồng bộ giữa 3 khớp này cho phép đầu bàn chân di chuyển đến bất kỳ tọa độ $(x, y, z)$ nào trong không gian làm việc của nó. Để đảm bảo sự đồng đều về lực, các động cơ được sử dụng thường có công suất định mức và mô-men xoắn tương đương nhau.

\section{Không gian làm việc }
Không gian làm việc chân robot là vùng không gian mà đầu bàn chân có thể vươn tới. Không gian làm việc của chân có thể được xác định bằng phương pháp phân tích hình học khi đã biết giới hạn góc quay của các động cơ và các khớp.

Phương pháp động học thuận rất hữu ích trong việc xác định không gian làm việc vì cơ cấu chân điển hình có cấu trúc tương tự như một tay máy nối tiếp. Đối với cơ cấu chân 3 bậc tự do với các giới hạn góc $\theta_{i,min} \leq \theta_i \leq \theta_{i,max}$ ($i = 1, 2, 3$), không gian làm việc được xác định bằng cách quét toàn bộ miền góc khớp cho phép và tính tọa độ đầu chân bằng FK:

\begin{equation}
    \mathbf{p}_{foot} = \text{FK}(\theta_1, \theta_2, \theta_3)
\end{equation}

Kết quả quét tạo ra một đám mây điểm 3D có dạng \textit{vòng cung} (annular sector), đặc trưng cho cơ cấu serial. Tầm với tối đa của chân đạt được khi các khâu duỗi thẳng ($\theta_3 = 0$):

\begin{equation}
    r_{max} = l_2 + l_3
\end{equation}

trong khi tầm với tối thiểu xảy ra khi các khâu gập hoàn toàn:

\begin{equation}
    r_{min} = |l_2 - l_3|
\end{equation}

Không gian làm việc có ý nghĩa quan trọng trong thiết kế: nó xác định phạm vi chiều cao đứng, chiều dài bước và biên độ chuyển động của robot. Vùng đạt được rộng nhất nằm trong mặt phẳng sagittal ($X$-$Z$), trong khi chiều ngang theo phương $Y$ bị giới hạn bởi biên độ khớp hông (hip - $\theta_1$).



\section{Các phương thức di chuyển}
Robot bốn chân có rất nhiều phương thức di chuyển khác nhau tùy thuộc vào tốc độ và đặc điểm địa hình. Dưới đây là 5 dáng đi (\textit{gait}) phổ biến nhất được áp dụng trong kỹ thuật điều khiển robot.

\textbf{\textit{Dáng đi Trot (Chạy kiệu)}} 

Đây là dáng đi phổ biến nhất và cũng là dáng đi trọng tâm được tác giả lựa chọn để thực hiện. Ở dáng đi này, hai chân chéo nhau (ví dụ: chân trước-trái và chân sau-phải) sẽ đồng thời nhấc lên và tiếp đất cùng lúc. \textit{Trot gait} giúp robot di chuyển nhanh và mượt mà, nhưng đòi hỏi bài toán điều khiển cân bằng động liên tục do đa giác chân đế thu hẹp lại thành một đường chéo.
    
\begin{figure}[H]
    \centering
    \includegraphics[width=0.6\textwidth]{gait_planner/trot_gait_table.png}
    \caption{Giản đồ pha chuyển động của dáng đi Trot}
    \label{fig3}
\end{figure}

\textbf{\textit{Dáng đi Walk (Đi bộ)}} 

Là dáng đi chậm và ổn định nhất. Tại bất kỳ thời điểm nào, robot luôn có ít nhất 3 chân tiếp đất để tạo thành một chân đế vững chắc, trong khi 1 chân được nhấc lên bước tới. Các pha (\textit{phase}) của 4 chân thường lệch nhau một góc $90^\circ$ (tương ứng $0^\circ, 90^\circ, 180^\circ, 270^\circ$).
    
\begin{figure}[H]
    \centering
    \includegraphics[width=0.6\textwidth]{gait_planner/walk_gait_table.png}
    \caption{Giản đồ pha chuyển động của dáng đi Walk}
    \label{fig4}
\end{figure}

\textbf{\textit{Dáng đi Pace (Đi kiệu song song)}} 

Hai chân cùng một bên (trái hoặc phải) được nhấc lên và tiến về phía trước cùng lúc. Đặc điểm của dáng đi này là trọng tâm khối lượng (CoM - \textit{Center of Mass}) của robot sẽ bị lắc lư ngang theo trục Y trong quá trình di chuyển, do đó hệ thống cần có cơ chế bù trừ thăng bằng tốt.
    
\begin{figure}[H]
    \centering
    \includegraphics[width=0.6\textwidth]{gait_planner/pace_gait_table.png}
    \caption{Giản đồ pha chuyển động của dáng đi Pace}
    \label{fig5}
\end{figure}

\textbf{\textit{Dáng đi Bound (Nhảy vọt)}} 

Hai chân trước đồng thời nhấc lên vươn tới, ngay sau đó hai chân sau cũng nhấc lên tạo ra chuyển động nảy bật về phía trước, mô phỏng chuyển động của các loài thú săn mồi.
    
\begin{figure}[H]
    \centering
    \includegraphics[width=0.6\textwidth]{gait_planner/bound_gait_table.png}
    \caption{Giản đồ pha chuyển động của dáng đi Bound}
    \label{fig6}
\end{figure}

\textbf{\textit{Dáng đi Gallop (Phi mã)}} 

Đây là dáng đi có tốc độ cao nhất của robot bốn chân. Quá trình di chuyển sẽ xuất hiện các pha bay hoàn toàn trên không (\textit{flight phase}), tức là không có chân nào chạm đất. Dáng đi này đòi hỏi các cơ cấu chấp hành có mô-men xoắn và tốc độ cực kỳ cao.
    
\begin{figure}[H]
    \centering
    \includegraphics[width=0.6\textwidth]{gait_planner/gallop_gait_table.png}
    \caption{Giản đồ pha chuyển động của dáng đi Gallop}
    \label{fig7}
\end{figure}

\section{Lý thuyết ổn định cho robot có chân}

Phân tích ổn định là yêu cầu thiết yếu trong thiết kế và vận hành robot bốn chân. Khác với robot bánh xe có diện tích tiếp xúc liên tục với mặt đất, robot có chân chỉ tiếp xúc tại các điểm rời rạc, và diện tích đỡ thay đổi liên tục theo thời gian khi các chân chuyển pha.

\subsection{Ổn định tĩnh}

Robot có chân được xem là ổn định tĩnh tại một thời điểm nếu hình chiếu của trọng tâm (Center of Mass -- CoM) xuống mặt đất nằm bên trong đa giác đỡ (support polygon) --- đa giác lồi được tạo bởi các điểm tiếp xúc giữa các chân đang chạm đất và mặt đất.

\textbf{\textit{Biên ổn định tĩnh}} 

(Static Stability Margin -- SSM) được định nghĩa là khoảng cách nhỏ nhất từ hình chiếu CoM đến các cạnh của đa giác đỡ \cite{gonzalez2018}:

\begin{equation}
    \text{SSM} = \min_{i} \, d(\mathbf{p}_{CoM}, \overline{P_i P_{i+1}})
\end{equation}

trong đó $\mathbf{p}_{CoM}$ là hình chiếu CoM lên mặt đất, $\overline{P_i P_{i+1}}$ là cạnh thứ $i$ của đa giác đỡ, và $d(\cdot, \cdot)$ là khoảng cách từ điểm đến đoạn thẳng. Nếu SSM $> 0$, robot ổn định tĩnh; nếu SSM $\leq 0$, robot sẽ đổ.

\textbf{\textit{Ứng dụng cho dáng đi Trot}} 

Trong dáng đi trot, tại mỗi thời điểm có hai chân chéo nhau chạm đất. Đa giác đỡ suy biến thành một đoạn thẳng nối hai bàn chân. Về mặt lý thuyết, điều này có nghĩa là robot đang ở trạng thái biên của ổn định tĩnh (SSM $\approx 0$ khi CoM nằm đúng trên đoạn thẳng). Trong thực tế, diện tích tiếp xúc hữu hạn của bàn chân và quán tính của thân robot tạo ra một vùng ổn định hẹp nhưng đủ để duy trì chuyển động.

\begin{figure}[H]
    \centering
    \begin{subfigure}[b]{0.48\textwidth}
        \centering
        \includegraphics[width=\textwidth]{kinematics/static_stability_a.png}
        \caption{Walk Gait -- 3 stance legs}
    \end{subfigure}
    \hfill
    \begin{subfigure}[b]{0.48\textwidth}
        \centering
        \includegraphics[width=\textwidth]{kinematics/static_stability_b.png}
        \caption{Trot Gait -- 2 diagonal stance legs}
    \end{subfigure}
    \caption{Minh họa Đa giác đỡ và Biên ổn định tĩnh (SSM) của robot bốn chân}
    \label{fig:static_stability}
\end{figure}

Trong nghiên cứu này, robot hoạt động ở tốc độ thấp với dáng đi Trot nên phân tích ổn định tĩnh là đủ cho mục đích thiết kế. Bộ điều khiển PID cân bằng kết hợp phản hồi IMU đóng vai trò bổ sung, giảm thiểu sai lệch hướng của thân robot khi di chuyển, chủ động đưa trọng tâm về gần đường nối hai bàn chân đang chạm đất.

\section{Lý thuyết bộ điều khiển PID}

Bộ điều khiển PID (Proportional--Integral--Derivative) là một trong những thuật toán điều khiển phản hồi được sử dụng rộng rãi và hiệu quả nhất trong kỹ thuật tự động hóa \cite{astrom2006}. Cấu trúc sơ đồ khối lý thuyết của bộ điều khiển PID được thể hiện trong Hình~\ref{fig:pid_theory_diagram}. Bộ điều khiển liên tục tính toán giá trị sai lệch $e(t) = r(t) - y(t)$ giữa giá trị đặt (setpoint) $r(t)$ và giá trị đo được (process variable) $y(t)$, sau đó áp dụng tín hiệu hiệu chỉnh dựa trên ba thành phần tỷ lệ, tích phân và vi phân.

\begin{figure}[H]
    \centering
    \includegraphics[width=0.8\textwidth]{architecture/Proportional-Integral-Derivative.png}
    \caption{Cấu trúc sơ đồ khối của bộ điều khiển PID}
    \label{fig:pid_theory_diagram}
\end{figure}

\subsection{Luật điều khiển PID chuẩn}

Tín hiệu điều khiển tổng hợp $u(t)$ trong miền thời gian liên tục được biểu diễn bằng phương trình:

\begin{equation}
    u(t) = K_P \cdot e(t) + K_I \int_{0}^{t} e(\tau)\,d\tau + K_D \frac{de(t)}{dt}
    \label{eq:pid_continuous}
\end{equation}

trong đó $K_P$, $K_I$, và $K_D$ lần lượt là hệ số khuếch đại tỷ lệ, hệ số tích phân và hệ số vi phân. Mỗi thành phần đóng một vai trò riêng biệt trong việc định hình đáp ứng của hệ thống. Trước hết, thành phần tỷ lệ (Proportional - P) tạo ra tín hiệu hiệu chỉnh tỷ lệ thuận với sai lệch hiện tại; hệ số $K_P$ càng lớn, hệ thống phản ứng càng nhanh nhưng đồng thời có xu hướng dao động và độ vọt lố (overshoot) cao. Thành phần P đơn thuần không thể triệt tiêu hoàn toàn sai lệch ở trạng thái xác lập. Tiếp theo, thành phần tích phân (Integral - I) tích lũy sai lệch theo thời gian để triệt tiêu hoàn toàn sai số ở trạng thái xác lập (steady-state error). Tuy nhiên, khi sai lệch tồn tại trong thời gian dài do hệ thống bị bão hòa cơ cấu chấp hành, giá trị tích phân tích lũy quá lớn sẽ gây ra hiện tượng bão hòa tích phân (\textit{integral windup}), khiến hệ thống bị vọt lố nghiêm trọng và phục hồi chậm. Cuối cùng, thành phần vi phân (Derivative - D) dự đoán xu hướng thay đổi của sai lệch dựa trên đạo hàm. Thành phần này hoạt động như một ``phanh'' hãm, giúp giảm vọt lố và cải thiện thời gian xác lập, mặc dù nhược điểm chính của khâu D là rất nhạy cảm với nhiễu tần số cao trong tín hiệu đo đạc từ cảm biến.

\subsection{Dạng rời rạc và triển khai số}

Trong các hệ thống vi điều khiển nhúng với chu kỳ lấy mẫu $\Delta t$, bộ điều khiển PID không thể tính toán các phép vi phân và tích phân liên tục. Do đó, phương trình \ref{eq:pid_continuous} được chuyển đổi thành dạng rời rạc (discrete-time):

\begin{equation}
    u_k = K_P \cdot e_k + K_I \sum_{j=0}^{k} e_j \cdot \Delta t + K_D \frac{e_k - e_{k-1}}{\Delta t}
    \label{eq:pid_discrete}
\end{equation}

trong đó chỉ số $k$ biểu thị bước thời gian hiện tại. Phương trình trên là dạng triển khai cơ bản nhất. Tuy nhiên, để đảm bảo hệ thống thực tế hoạt động ổn định và tin cậy, các kỹ thuật cải tiến thường được áp dụng thêm.

\subsection{Các cải tiến PID trong hệ thống thực}

Trong điều kiện vận hành thực tế, đặc biệt đối với robot bốn chân chịu ảnh hưởng lớn từ nhiễu động học và giới hạn phần cứng, bộ điều khiển PID cơ bản (phương trình \ref{eq:pid_discrete}) không thể đáp ứng tốt. Dựa trên mô hình MATLAB/Simulink và thực nghiệm, chuỗi điều khiển PID trong nghiên cứu này được tích hợp các cơ chế cải tiến sau:

\textbf{\textit{Bộ lọc trung bình trượt (EMA - Exponential Moving Average)}}

Trong quá trình di chuyển (chế độ GAIT), thân robot luôn dao động tuần hoàn. Nếu bộ điều khiển lấy mốc thăng bằng tuyệt đối (0 rad), nó sẽ liên tục sinh ra tín hiệu chống lại chính nhịp bước tự nhiên của robot. Việc áp dụng bộ lọc EMA tạo ra một đường cơ sở động (dynamic baseline):
\begin{equation}
    \text{Baseline}_k = \alpha_{ema} \cdot \text{Baseline}_{k-1} + (1 - \alpha_{ema}) \cdot y_k
\end{equation}
Sai lệch đưa vào PID sẽ là $e_k = \text{Baseline}_k - y_k$, giúp bộ điều khiển bỏ qua các dao động tự nhiên và chỉ phản ứng khi thân robot thực sự mất trọng tâm.

\textbf{\textit{Vùng chết (Deadband)}}

Cảm biến IMU luôn tồn tại nhiễu nhỏ ngay cả khi hệ thống ở vị trí cân bằng. Vùng chết triệt tiêu các sai lệch không đáng kể này để tránh động cơ phải bù trừ liên tục (chatter), giảm mài mòn cơ khí và tiết kiệm năng lượng:
\begin{equation}
    e_{dz} = \begin{cases}
        e - d & \text{nếu } e > d \\
        0 & \text{nếu } |e| \leq d \\
        e + d & \text{nếu } e < -d
    \end{cases}
\end{equation}
trong đó $d$ là bán kính của vùng chết.

\textbf{\textit{Chống bão hòa (Anti-windup) và Rò rỉ tích phân (Integral Leakage)}}

Để ngăn khâu tích phân (I) giữ lại các sai số cũ quá lâu hoặc bị bão hòa khi đạt giới hạn cơ cấu chấp hành, hệ số rò rỉ $\lambda \in (0, 1)$ được áp dụng để làm giảm dần giá trị tích phân theo thời gian (hàm mũ):
\begin{equation}
    \mathcal{I}_k = \lambda \cdot \mathcal{I}_{k-1} + e_k \cdot \Delta t
\end{equation}
Đồng thời, giá trị $\mathcal{I}_k$ được khống chế cứng (clamping) trong khoảng $[-\mathcal{I}_{max}, \mathcal{I}_{max}]$ nhằm đảm bảo an toàn.

\textbf{\textit{Reset tích phân tại giao điểm không (Zero-Crossing Reset)}}

Trong chế độ di chuyển tuần hoàn như dáng đi Trot, thân robot lắc qua lại liên tục qua trục cân bằng. Để tránh khâu tích phân gây ra quá điều chỉnh (overshoot) khi đổi chiều, hệ thống sẽ tự động xóa giá trị tích lũy ($\mathcal{I}_k = 0$) ngay khi sai số $e_k$ đổi dấu (vượt qua giao điểm không).

\textbf{\textit{Bão hòa đầu ra và Lọc thông thấp (Saturation \& Output LPF)}}

Tín hiệu hiệu chỉnh từ PID có thể chứa các gai nhiễu (chủ yếu từ khâu D). Trước khi xuất xuống hệ thống động học, tín hiệu được giới hạn biên độ chặt chẽ ($[-u_{max}, u_{max}]$) để không phá vỡ quỹ đạo sải bước nguyên vẹn. Sau đó, tín hiệu tiếp tục đi qua bộ lọc thông thấp (Low-Pass Filter) để làm mượt, đảm bảo động cơ vận hành êm ái, tránh hiện tượng trượt chân hoặc hụt bước.

\section{Cảm biến quán tính IMU và biểu diễn hướng}

\subsection{Nguyên lý hoạt động}

Bộ đo lường quán tính (Inertial Measurement Unit -- IMU) là thiết bị kết hợp nhiều loại cảm biến để đo gia tốc, vận tốc góc và hướng của vật thể trong không gian ba chiều. Một IMU tiêu chuẩn 9 trục thường bao gồm ba thành phần chính \cite{madgwick2011}. Thứ nhất là gia tốc kế (Accelerometer) dùng để đo gia tốc tuyến tính theo ba trục ($a_x, a_y, a_z$), bao gồm cả thành phần gia tốc trọng trường; từ các phép đo này, góc nghiêng tĩnh (roll, pitch) có thể được suy ra thông qua hướng của vector trọng lực. Thứ hai là con quay hồi chuyển (Gyroscope) làm nhiệm vụ đo vận tốc góc ($\omega_x, \omega_y, \omega_z$) quanh ba trục; mặc dù việc tích phân vận tốc góc theo thời gian cho ra góc quay, do sai số tích lũy (drift), phương pháp này không đáng tin cậy trong thời gian dài. Cuối cùng là từ kế (Magnetometer) có chức năng đo cường độ từ trường Trái Đất theo ba trục, cung cấp thông tin về hướng Bắc từ, từ đó cho phép hệ thống xác định góc hướng (yaw).

\subsection{Biểu diễn hướng: Quaternion và góc Euler}

Hướng của một vật rắn trong không gian ba chiều có thể được biểu diễn bằng nhiều phương pháp, trong đó phổ biến nhất là góc Euler và quaternion.

\textbf{\textit{Góc Euler}} 

Sử dụng ba góc quay liên tiếp quanh các trục cố định hoặc trục động để mô tả hướng. Trong hệ quy ước ZYX (yaw--pitch--roll), ma trận xoay tổng hợp là:

\begin{equation}
    R = R_z(\psi) \cdot R_y(\theta) \cdot R_x(\phi)
\end{equation}

trong đó $\phi$ là góc lăn (roll), $\theta$ là góc nghiêng (pitch), và $\psi$ là góc hướng (yaw). Nhược điểm chính của biểu diễn Euler là hiện tượng khóa gimbal (gimbal lock), xảy ra khi $\theta = \pm 90^\circ$, khiến hai trục quay trùng nhau và mất một bậc tự do.

\textbf{\textit{Quaternion}} 

$\mathbf{q} = (w, x, y, z)$ với $|\mathbf{q}| = 1$ biểu diễn phép quay một góc $\alpha$ quanh trục đơn vị $\hat{n} = (n_x, n_y, n_z)$:

\begin{equation}
    \mathbf{q} = \left(\cos\frac{\alpha}{2},\; n_x \sin\frac{\alpha}{2},\; n_y \sin\frac{\alpha}{2},\; n_z \sin\frac{\alpha}{2}\right)
\end{equation}

Quaternion không bị hiện tượng khóa gimbal, tính toán nội suy mượt mà hơn, và yêu cầu ít phép toán hơn so với ma trận xoay. Công thức chuyển đổi từ quaternion sang góc Euler thường dùng trong robotics là:

\begin{align}
    \phi   &= \text{atan2}\left(2(w x + y z),\; 1 - 2(x^2 + y^2)\right) \\
    \theta &= \arcsin\left(2(w y - z x)\right) \\
    \psi   &= \text{atan2}\left(2(w z + x y),\; 1 - 2(y^2 + z^2)\right)
\end{align}

\subsection{Bộ lọc cảm biến}

Dữ liệu thô từ IMU chứa nhiễu tần số cao và sai số trôi (drift), do đó cần được xử lý qua bộ lọc trước khi sử dụng. Có nhiều phương pháp phổ biến được ứng dụng trong thực tế. Khái niệm cơ bản nhất là bộ lọc bổ sung (Complementary Filter), kết hợp dữ liệu gia tốc kế (tin cậy ở tần số thấp) và con quay hồi chuyển (tin cậy ở tần số cao) thông qua hệ số trọng số $\alpha$:
\begin{equation}
    \hat{\theta}_k = \alpha \cdot (\hat{\theta}_{k-1} + \omega_k \Delta t) + (1-\alpha) \cdot \theta_{acc,k}
\end{equation}
Bên cạnh đó, bộ lọc Kalman là một phương pháp ước lượng tối ưu dựa trên mô hình trạng thái, kết hợp dự đoán (prediction) từ mô hình động học với cập nhật (update) từ phép đo. Trong đó, bộ lọc Kalman mở rộng (EKF) thường được sử dụng cho các hệ phi tuyến. Một giải pháp thay thế hiệu quả là bộ lọc Madgwick \cite{madgwick2011}, một thuật toán tính toán nhẹ sử dụng quy hoạch gradient (gradient descent) để tối ưu hóa quaternion ước lượng, rất phù hợp với các hệ thống nhúng có tài nguyên hạn chế. Hiện nay, giải pháp tối ưu hóa phần cứng là sử dụng Sensor Fusion tích hợp (On-chip); một số module IMU cao cấp như BNO055 được tích hợp sẵn bộ xử lý sensor fusion để tính toán và cung cấp trực tiếp quaternion hoặc góc Euler đã được lọc nhiễu, giúp giảm đáng kể khối lượng tính toán cho vi xử lý chính.

\section{Tổng quan về nền tảng ROS~2 và Gazebo}

ROS~2 (Robot Operating System 2) \cite{ros2docs} là nền tảng phần mềm mã nguồn mở cung cấp kiến trúc truyền thông phân tán thời gian thực dựa trên middleware DDS. Việc sử dụng ROS~2 cho phép chia nhỏ các tác vụ điều khiển robot thành các Node độc lập (ví dụ: Node tính toán động học, Node điều khiển cân bằng) giao tiếp lỏng ghép với nhau thông qua cơ chế Publish/Subscribe, qua đó tăng cường độ tin cậy và khả năng mở rộng của hệ thống.

Song song với ROS~2, Gazebo \cite{gazebo} được sử dụng làm môi trường mô phỏng vật lý 3D. Nhờ khả năng tích hợp chặt chẽ với ROS~2 và hỗ trợ mô phỏng động lực học chính xác (trọng lực, ma sát, va chạm), đồ án áp dụng quy trình phát triển ``sim-to-real''. Thuật toán điều khiển và mô hình URDF của robot được kiểm thử toàn diện trên Gazebo để loại trừ các lỗi logic trước khi triển khai xuống máy tính nhúng và cơ cấu chấp hành thực tế, giúp tiết kiệm thời gian và đảm bảo an toàn phần cứng.


Mục tiêu của bài toán động học thuận của robot là xác định biểu diễn vị trí của cơ cấu chấp hành cuối. Trong nghiên cứu này, phương pháp biến đổi tọa độ Modified Denavit--Hartenberg (MD-H) \cite{denavit1955, craig2005} được áp dụng để giải quyết bài toán này.
\section{Mô hình động học}
\subsection{Thiết lập hệ tọa độ MD-H}

Nhóm xây dựng một mô hình toán học và cơ khí của robot bốn chân phân bố trong không gian Descartes. Cấu trúc vật lý của robot bao gồm kích thước bệ đỡ với chiều dài $L=209$ mm, chiều rộng $W=191$ mm, và chiều dài các khâu chân lần lượt là $L_1=26$ mm (Coxa/Hip), $L_2=106$ mm (Femur/Thigh), và $L_3=125$ mm (Tibia/Calf).

\begin{figure}[H]
    \centering
    \includegraphics[width=0.7\textwidth]{kinematics/model_reference_frame.png}
    \caption{Mô hình mô phỏng với các hệ tọa độ tham chiếu được gắn tại từng khớp của robot.}
    \label{fig20}
\end{figure}

Theo quy ước chuẩn của hệ điều hành ROS (Robot Operating System), các trục tọa độ trong Hình \ref{fig20} được biểu diễn bằng màu sắc: Trục $X$ (Đỏ) chỉ hướng tiến về phía trước, Trục $Y$ (Xanh lá) chỉ hướng sang trái, và Trục $Z$ (Xanh dương) chỉ hướng lên trên. Các hệ tọa độ cấu thành nên mô hình động học MD-H của chân được gắn cụ thể trên hình, bao gồm: Hệ tọa độ thân $\{B\}$, cụm khớp hông (Roll và Pitch) đại diện bởi $\{0\}, \{1\}, \{2\}$, khớp gối $\{3\}$ và mũi bàn chân $\{4\}$.

Dựa theo hướng tiến về phía trước, chân trước bên trái (LF), chân trước bên phải (RF), chân sau bên trái (LB) và chân sau bên phải (RB) lần lượt được đánh số là 1, 2, 3 và 4. Hệ tọa độ thân $\{B\}$ (hoặc $\{O_B\}$) được thiết lập tại tâm hình học của thân, và hệ tọa độ gốc của chân $\{0\}$ (hoặc $\{O_{i,0}\}$) được thiết lập tại khớp nối giữa khớp hông và thân, trong đó $i = 1,2,3,4$.

Trong phương pháp biến đổi tọa độ Modified Denavit-Hartenberg (MD-H) \cite{craig2005}, chuỗi động học được mô tả một cách toán học thông qua các phép biến đổi tọa độ chính xác. Mỗi khâu liên kết được xác định bởi bốn tham số như trong Hình~\ref{fig21}, theo định nghĩa chuẩn của MD-H, bao gồm: $a_{i-1}$ là khoảng cách từ trục $Z_{i-1}$ đến trục $Z_i$ đo dọc theo trục $X_{i-1}$; $\alpha_{i-1}$ là góc quay từ trục $Z_{i-1}$ đến trục $Z_i$ đo quanh trục $X_{i-1}$; $d_i$ là khoảng cách từ trục $X_{i-1}$ đến trục $X_i$ đo dọc theo trục $Z_i$; và cuối cùng, $\theta_i$ là góc quay từ trục $X_{i-1}$ đến trục $X_i$ đo quanh trục $Z_i$.

\begin{figure}[H]
    \centering
    \includegraphics[width=0.6\textwidth]{kinematics/dh_coordinate_system.png }
    \caption{Mô hình một chân của robot bốn chân được xây dựng bằng phương pháp MD-H.}
    \label{fig21}
\end{figure}

Các tham số MD-H của từng khâu liên kết cho một chân robot được thiết lập và trình bày chi tiết trong Bảng \ref{tab:3}, phản ánh chính xác sự khác biệt về cấu trúc động học giữa các chân bên trái và bên phải (tham số $d_2$ mang dấu đối nhau do hệ quả của phép đối xứng gương).

\begin{table}[H]
\centering
\caption{Tham số Modified Denavit-Hartenberg (MD-H)}
\label{tab:3}
\renewcommand{\arraystretch}{1.4}
\begin{tabular}{l c c c c c}
\toprule
\textbf{Link ($i$)} & $\boldsymbol{\alpha_{i-1}}$ & $\boldsymbol{a_{i-1}}$ & \multicolumn{2}{c}{$\boldsymbol{d_i}$} & $\boldsymbol{\theta_i}$ \\
\cmidrule(lr){4-5}
& & & \textbf{Left Limbs} & \textbf{Right Limbs} & \\
\midrule
1 (Coxa)  & $0^\circ$ & $0$   & $0$    & $0$   & $\theta_1^*$ \\
2 (Femur) & $90^\circ$& $0$   & $L_1$  & $-L_1$& $\theta_2^*$ \\
3 (Tibia) & $0^\circ$ & $L_2$ & $0$    & $0$   & $\theta_3^*$ \\
4 (Foot)  & $0^\circ$ & $L_3$ & $0$    & $0$   & $0$ \\
\bottomrule
\multicolumn{6}{l}{\small $^*$ Variable articulation angle} \\
\end{tabular}
\end{table}

\subsection{Giải động học thuận}

Khi giải các phương trình động học thuận bằng phương pháp MD-H, chúng thường được 
phân rã thành các bài toán con. Phép biến đổi từ hệ tọa độ khâu $\{i-1\}$ sang hệ 
tọa độ $\{i\}$ tạo ra một ma trận biến đổi tọa độ chỉ liên quan đến bốn tham số 
liên kết $^{i-1}_{i}T$. Dựa trên định nghĩa của từng tham số liên kết ở trên, ta có:

\begin{equation}
^{i-1}_{i}T = Trans_x\{a_{i-1}\} \, Rot_x\{\alpha_{i-1}\} \, Trans_z\{d_i\} \, Rot_z\{\theta_i\}
\end{equation}
\begin{equation}
{}^{i-1}_{i}T=
\begin{bmatrix}
\cos\theta_i & -\sin\theta_i & 0 & a_{i-1} \\
\sin\theta_i\cos\alpha_{i-1} & \cos\theta_i\cos\alpha_{i-1} & -\sin\alpha_{i-1} & -d_i\sin\alpha_{i-1} \\
\sin\theta_i\sin\alpha_{i-1} & \cos\theta_i\sin\alpha_{i-1} & \cos\alpha_{i-1} & d_i\cos\alpha_{i-1} \\
0 & 0 & 0 & 1
\end{bmatrix}
\end{equation}
Áp dụng cho cấu trúc của một chân robot, các ma trận biến đổi thành phần được xác định như sau:

\begin{equation}
{}^{0}_{1}T=
\begin{bmatrix}
\cos\theta_1 & -\sin\theta_1 & 0 & 0 \\
\sin\theta_1 & \cos\theta_1 & 0 & 0 \\
0 & 0 & 1 & 0 \\
0 & 0 & 0 & 1
\end{bmatrix}
\end{equation}

\begin{equation}
{}^{1}_{2}T=
\begin{bmatrix}
\cos\theta_2 & -\sin\theta_2 & 0 & 0 \\
0 & 0 & -1 & -L_1 \\
\sin\theta_2 & \cos\theta_2 & 0 & 0 \\
0 & 0 & 0 & 1
\end{bmatrix}
\end{equation}

\begin{equation}
{}^{2}_{3}T=
\begin{bmatrix}
\cos\theta_3 & -\sin\theta_3 & 0 & L_2 \\
\sin\theta_3 & \cos\theta_3 & 0 & 0 \\
0 & 0 & 1 & 0 \\
0 & 0 & 0 & 1
\end{bmatrix}
\end{equation}

\begin{equation}
{}^{3}_{4}T=
\begin{bmatrix}
1 & 0 & 0 & L_3 \\
0 & 1 & 0 & 0 \\
0 & 0 & 1 & 0 \\
0 & 0 & 0 & 1
\end{bmatrix}
\end{equation}
Theo Hình \ref{fig20}, ma trận biến đổi từ hệ tọa độ $\{O_{i,0}\}$ đến hệ tọa độ thân $\{O_B\}$ được viết như sau:

\begin{equation}
^{B}_{0}T=
\begin{bmatrix}
0 & 0 & 1 & \frac{L}{2} \\
0 & -1 & 0 & \frac{W}{2} \\
1 & 0 & 0 & 0 \\
0 & 0 & 0 & 1
\end{bmatrix}
\end{equation}

Do đó, ma trận biến đổi của hệ tọa độ đầu chân so với hệ tọa độ thân $\{O_B\}$ được xác định như sau:

\begin{equation}
^{B}_{4}T = {}^{B}_{0}T \, {}^{0}_{4}T
\end{equation}

\begin{equation}
{}^{B}_{4}T =
\begin{bmatrix}
S_{23} & c_{23} & 0 & \frac{L}{2} + L_2 S_2 + L_3 S_{23} \\
- S_1 c_{23} & S_1 S_{23} & c_1 & \frac{W}{2} + L_1 c_1 - L_2 S_1 c_2 - L_3 S_1 c_{23} \\
c_1 c_{23} & - c_1 S_{23} & S_1 & L_1 S_1 + L_2 c_1 c_2 + L_3 c_1 c_{23} \\
0 & 0 & 0 & 1
\end{bmatrix}
\end{equation}

\begin{equation}
{}^{B}_{4}T=
\begin{bmatrix}
R & P \\
0 & 1
\end{bmatrix}
\end{equation}

\begin{equation}
{}^{B}_{4}T=
\begin{bmatrix}
 &  &  & P_x \\
 & R &  & P_y \\
 &  &  & P_z \\
0 & 0 & 0 & 1
\end{bmatrix}
\end{equation}

\begin{equation}
R=
\begin{bmatrix}
S_{23} & c_{23} & 0 \\
- S_1 c_{23} & S_1 S_{23} & c_1 \\
c_1 c_{23} & - c_1 S_{23} & S_1
\end{bmatrix}
\end{equation}

\begin{equation}
P =
\begin{bmatrix}
P_x^B \\
P_y^B \\
P_z^B
\end{bmatrix}
=
\begin{bmatrix}
\dfrac{L}{2}+L_2 s_2 + L_3 s_{23} \\
\dfrac{W}{2}+L_1 c_1 - L_2 s_1 c_2 - L_3 s_1 c_{23} \\
L_1 s_1 + L_2 c_1 c_2 + L_3 c_1 c_{23}
\end{bmatrix}
\end{equation}

\begin{equation}
P_{LF}=
\begin{bmatrix}
P_x \\
P_y \\
P_z
\end{bmatrix}
=
\begin{bmatrix}
L_1\sin\theta_1 + L_2\cos\theta_1\cos\theta_2 + L_3\cos\theta_1\cos(\theta_2+\theta_3) \\
- L_1\cos\theta_1 + L_2\sin\theta_1\cos\theta_2 + L_3\sin\theta_1\cos(\theta_2+\theta_3) \\
L_2\sin\theta_2 + L_3\sin(\theta_2+\theta_3)
\end{bmatrix}
\label{eq:1}
\end{equation}

\begin{equation}
P_{RF}=
\begin{bmatrix}
P_x \\
P_y \\
P_z
\end{bmatrix}
=
\begin{bmatrix}
- L_1\sin\theta_1 + L_2\cos\theta_1\cos\theta_2 + L_3\cos\theta_1\cos(\theta_2+\theta_3) \\
L_1\cos\theta_1 + L_2\sin\theta_1\cos\theta_2 + L_3\sin\theta_1\cos(\theta_2+\theta_3) \\
L_2\sin\theta_2 + L_3\sin(\theta_2+\theta_3)
\end{bmatrix}
\label{eq:2}
\end{equation}

\begin{equation}
P_{LB}=
\begin{bmatrix}
P_x \\
P_y \\
P_z
\end{bmatrix}
=
\begin{bmatrix}
L_1\sin\theta_1 + L_2\cos\theta_1\cos\theta_2 + L_3\cos\theta_1\cos(\theta_2+\theta_3) \\
- L_1\cos\theta_1 + L_2\sin\theta_1\cos\theta_2 + L_3\sin\theta_1\cos(\theta_2+\theta_3) \\
L_2\sin\theta_2 + L_3\sin(\theta_2+\theta_3)
\end{bmatrix}
\label{eq:3}
\end{equation}

\begin{equation}
P_{RB}=
\begin{bmatrix}
P_x \\
P_y \\
P_z
\end{bmatrix}
=
\begin{bmatrix}
- L_1\sin\theta_1 + L_2\cos\theta_1\cos\theta_2 + L_3\cos\theta_1\cos(\theta_2+\theta_3) \\
L_1\cos\theta_1 + L_2\sin\theta_1\cos\theta_2 + L_3\sin\theta_1\cos(\theta_2+\theta_3) \\
L_2\sin\theta_2 + L_3\sin(\theta_2+\theta_3)
\end{bmatrix}
\label{eq:4}
\end{equation}

Các ma trận trên biểu diễn vị trí mũi bàn chân của cả 4 chân trong hệ tọa độ cục bộ của từng khớp hông. Do tính chất đối xứng, phương trình động học của chân sau bên trái giống với chân trước bên trái, và chân sau bên phải giống với chân trước bên phải.

\section{Động học nghịch }

Bài toán động học nghịch (Inverse Kinematics -- IK) là bài toán xác định các giá trị góc khớp $(\theta_1, \theta_2, \theta_3)$ khi biết tọa độ đầu bàn chân $(x, y, z)$ trong hệ tọa độ thân robot. Đây là bài toán ngược của động học thuận, và thường có nhiều nghiệm (do tính phi tuyến của các phương trình lượng giác) hoặc không có nghiệm (khi điểm đích nằm ngoài không gian làm việc).

Trong nghiên cứu này, phương pháp giải tích (analytical method) được sử dụng để giải bài toán IK. Phương pháp này tận dụng cấu trúc chuỗi hở của chân robot để giải trực tiếp hệ phương trình lượng giác, từ đó thu được kết quả nhanh chóng và chính xác dạng đóng. Các bước tính toán chi tiết sẽ được trình bày cụ thể dưới đây.

Khi thiết kế phương trình động học nghịch cho robot bốn chân, một điểm đặc biệt cần lưu ý là các chân trái và chân phải được lắp ráp đối xứng nhau như hình ảnh qua gương. Do tính chất này, ta không thể áp dụng chung một công thức toán học y hệt nhau cho cả 4 chân. Cụ thể, có ba yếu tố vật lý tạo ra sự khác biệt. Thứ nhất là vị trí gắn chân trên thân; do mỗi chân được gắn ở một góc khác nhau, khoảng cách từ tâm robot đến các gốc chân có sự thay đổi về dấu, tương ứng với một nửa chiều dài ($\pm L/2$) và một nửa chiều rộng ($\pm W/2$). Thứ hai là cấu hình của khớp hông (Coxa); do lắp đặt đối xứng, với cùng một góc quay, chân trái và chân phải sẽ vung về hai hướng ngược nhau, do đó để công thức toán học tính đúng hướng, tham số chiều dài khâu hông $L_1$ sẽ mang dấu âm ($-L_1$) đối với các chân bên trái và mang dấu dương ($+L_1$) đối với các chân bên phải. Cuối cùng là cấu hình gập gối (Tibia); để duy trì hình dáng gập gối đồng nhất cho cả hai bên (để gối không bị bẻ ngược), toán hạng $\sin\theta_3$ bắt buộc phải nhận hai dấu trái ngược nhau giữa chân trái ($-$) và chân phải ($+$).

Thay vì phải xây dựng 4 hệ phương trình động học hoàn toàn độc lập, giải pháp tối ưu là thiết kế một bộ công thức tổng quát duy nhất. Sau đó, tùy thuộc vào việc đang tính toán cho chân nào, thuật toán chỉ việc tra bảng để đưa các hệ số dấu tương ứng vào công thức nhằm tìm ra góc khớp chính xác.

Cụ thể trong triển khai thực tế, hàm IK nhận tọa độ đầu bàn chân $(x_B, y_B, z_B)$ trong hệ tọa độ thân robot $\{O_B\}$. Dựa vào tính chất đối xứng trên, tọa độ này sẽ được biến đổi sang hệ tọa độ cục bộ của chân $(x', y', z')$ bằng cách cộng trừ các offset vị trí gốc chân tương ứng và hoán đổi trục cho phù hợp với quy ước MD-H. Phương trình biến đổi cho mỗi chân được viết như sau:

\begin{equation}
x' = z_B, \qquad
y' = \sigma_W \cdot \frac{W}{2} - y_B, \qquad
z' = \sigma_L \cdot \frac{L}{2} + x_B
\label{eq:body_to_leg}
\end{equation}

trong đó $\sigma_W$ và $\sigma_L$ là các hệ số dấu phụ thuộc vào vị trí chân, được trình bày trong Bảng~\ref{tab:ik_signs}.

\begin{table}[h]
\centering
\caption{Hệ số dấu biến đổi tọa độ và dấu IK cho bốn chân}
\label{tab:ik_signs}
\renewcommand{\arraystretch}{1.3}
\begin{tabular}{|l|c|c|c|c|}
\hline
\textbf{Chân} & $\sigma_W$ & $\sigma_L$ & Dấu $L_1$ trong $\theta_1$ & Dấu $\sin\theta_3$ \\ \hline
Left-Front  & $+1$ & $-1$ & $-L_1$ & $-1$ \\ \hline
Left-Behind & $+1$ & $+1$ & $-L_1$ & $-1$ \\ \hline
Right-Front & $-1$ & $-1$ & $+L_1$ & $+1$ \\ \hline
Right-Behind & $-1$ & $+1$ & $+L_1$ & $+1$ \\ \hline
\end{tabular}
\end{table}

\textbf{\textit{Quá trình giải hệ phương trình}} 

Sau khi biến đổi tọa độ, bài toán IK được giải thống nhất cho cả bốn chân thông qua bốn bước sau. Các hệ số dấu tham chiếu Bảng~\ref{tab:ik_signs}.

\textbf{\textit{Bước 1 -- Giải $\theta_1$}} 

Từ phép chiếu tọa độ $(x', y')$ lên mặt phẳng vuông góc với trục hip:
\begin{equation}
\theta_1 =
\mathrm{atan2}(-x',\, y')
+
\mathrm{atan2}
\left(
-\sqrt{x'^2 + y'^2 - L_1^2},\; \pm L_1
\right)
\label{eq:ik_theta1}
\end{equation}
trong đó dấu $\pm L_1$ là $-L_1$ cho chân trái và $+L_1$ cho chân phải.

\textbf{\textit{Bước 2 -- Tính biến trung gian}}
\begin{equation}
\begin{aligned}
p_1 &= x' \cos\theta_1 + y' \sin\theta_1 \\
p_2 &= \sigma_{p2} \cdot z'
\end{aligned}
\label{eq:ik_p1p2}
\end{equation}
trong đó $\sigma_{p2} = +1$ cho chân trái và $\sigma_{p2} = -1$ cho chân phải.

\textbf{\textit{Bước 3 -- Giải $\theta_3$}} 

Áp dụng định lý hàm số cosine:
\begin{equation}
\cos\theta_3 =
\frac{p_1^2 + p_2^2 - L_2^2 - L_3^2}{2 L_2 L_3}
\end{equation}

\begin{equation}
\sin\theta_3 =
\sigma_{s3} \sqrt{1 - \cos^2\theta_3}
\end{equation}

\begin{equation}
\theta_3 = \mathrm{atan2}(\sin\theta_3,\, \cos\theta_3)
\end{equation}
trong đó $\sigma_{s3} = -1$ cho chân trái và $\sigma_{s3} = +1$ cho chân phải, nhằm duy trì cấu hình gối đối xứng.

\textbf{\textit{Bước 4 -- Giải $\theta_2$}}
\begin{equation}
\theta_2 =
\mathrm{atan2}(p_2,\, p_1)
-
\mathrm{atan2}(L_3 \sin\theta_3,\, L_2 + L_3 \cos\theta_3)
\label{eq:ik_theta2}
\end{equation}

\section{Quy hoạch dáng đi và phân tích chuyển động }

Quy hoạch dáng đi đóng vai trò then chốt đối với sự di chuyển ổn định của robot bốn chân. Trong tự nhiên, các loài động vật có vú bốn chân lựa chọn các dáng đi khác nhau tùy thuộc vào vận tốc di chuyển mong muốn và độ gồ ghề của bề mặt địa hình nhằm giảm thiểu năng lượng tiêu thụ mà vẫn đạt được mục tiêu di chuyển. Trên cơ sở đó, các mẫu dáng đi phổ biến của động vật đã được trích xuất để làm nền tảng cho việc quy hoạch dáng đi trên robot bốn chân.

\subsection{Lựa chọn dáng đi Trot}
Dựa trên các phân tích lý thuyết về các phương thức di chuyển đã trình bày ở Chương 2, tác giả lựa chọn dáng đi \textbf{trot} làm trọng tâm nghiên cứu \cite{marhefka1999}. Đây là kiểu dáng đi ``hai nhịp'' điển hình, trong đó hai chân chéo nhau có đặc tính chuyển động giống hệt nhau, với chân trước bên trái (LF) và chân sau bên phải (RB) hoạt động đồng nhất, còn chân trước bên phải (RF) và chân sau bên trái (LB) hoạt động đồng nhất, như được thể hiện trong Hình~\ref{fig22}.

Về mặt lý thuyết, hệ số chu kỳ nhiệm vụ (duty cycle coefficient) trong dáng đi này không chịu bất kỳ hạn chế nào, và thỏa mãn điều kiện:
\begin{equation}
    0 < \beta < 1
\end{equation}
Do đó, robot có dải tốc độ rộng và có thể lựa chọn các tốc độ khác nhau cho các địa hình khác nhau. Ngoài ra, các thực nghiệm về lượng oxy tiêu thụ của ngựa ở các dáng đi khác nhau (walk, trot và gallop) cho thấy dáng đi trot có mức tiêu thụ năng lượng tương đối thấp trong một dải tốc độ khá lớn. Điều này chứng minh rằng dáng đi này có hiệu suất năng lượng cao về mặt lý thuyết.

\begin{figure}[H]
    \centering
    \includegraphics[width=0.6\textwidth]{gait_planner/trot_gait_timing.png}
    \caption{Sơ đồ thời gian của dáng đi trot với $\beta = 0.5$, trong đó $\beta = \frac{T_{st}}{T}$ và $T_{st}$ là thời gian của pha chống đỡ.}
    \label{fig22}
\end{figure}

Tùy thuộc vào giá trị của hệ số chu kỳ làm việc, dáng đi trot chéo 
có thể được chia thành ba loại.

Khi $\beta < 0.5$, robot xuất hiện một pha bay ngắn, tức là thời điểm 
mà cả bốn chân đều ở pha vung. Khi đó tốc độ di chuyển của robot tăng 
đáng kể, tuy nhiên sự xuất hiện của pha bay này cũng làm tăng đáng kể độ khó trong 
việc điều khiển.

Khi $\beta > 0.5$, robot có bốn giai đoạn chống đỡ khi chân vung chuyển 
đổi trạng thái. Điều này giúp tăng độ ổn định nhưng đồng thời làm giảm tốc độ di 
chuyển của robot.

Khi $\beta = 0.5$, hai cặp chân chéo luân phiên chuyển đổi liên tục giữa pha vung 
và pha chống đỡ. Trạng thái này nằm giữa hai trường hợp trên và đảm bảo sự cân 
bằng giữa tốc độ di chuyển và độ ổn định của robot, do đó đây là dáng đi được 
sử dụng phổ biến nhất.

\section{Quy hoạch quỹ đạo cho robot bốn chân}

Chuyển động của chân robot về cơ bản được mô tả thông qua quỹ đạo chuyển động của robot. 
Trong đó, quỹ đạo là sự biến thiên theo thời gian của vị trí, vận tốc và gia tốc 
của từng bậc tự do của robot.

Khi quá trình mô phỏng được thực hiện mà không có một phương pháp quy hoạch quỹ đạo phù hợp, 
chuyển động của robot thường trở nên ngẫu nhiên và không đều, đồng thời robot sẽ khó duy trì 
được sự cân bằng. Nguyên nhân của hiện tượng này là do sự thay đổi đột ngột của vận tốc trong 
không gian Descartes.

\begin{figure}[H]
    \centering
    \includegraphics[width=0.6\textwidth]{gait_planner/single_leg_trajectory.png }
    \caption{Phân tích quỹ đạo của một chân}
    \label{fig23}
\end{figure}

Việc quy hoạch quỹ đạo là cần thiết cho quá trình phân tích chuyển động dựa trên 
động học. Do đó, quỹ đạo được tạo ra với sự hỗ trợ của các điểm trung gian nhằm 
duy trì dạng chuyển động mượt mà. Vì robot bốn chân có 3 bậc tự do nên chuyển 
động của robot có sự thay đổi trên cả ba trục trong không gian Descartes.

\subsection{Nội suy đa thức bậc ba}

Trong quá trình xây dựng quỹ đạo chuyển động, phương pháp nội suy đa thức thường được sử dụng để tạo ra đường cong chuyển 
động liên tục giữa các điểm. Một trong những phương pháp phổ biến là nội suy 
đa thức bậc ba. Dạng tổng quát của đa thức bậc ba được 
biểu diễn như sau:

\[
f(t) = a_3 t^3 + a_2 t^2 + a_1 t + a_0
\]

Trong đó $a_0, a_1, a_2, a_3$ là các hệ số được xác định dựa trên bốn điều kiện biên 
gồm vị trí và vận tốc tại các điểm đầu và cuối của quỹ đạo:

\begin{align}
x(0) &= x_0, \quad x(T) = x_f \\
\dot{x}(0) &= 0, \quad \dot{x}(T) = 0
\end{align}

Phương pháp này giúp tạo ra quỹ đạo liên tục và đảm bảo tính liên tục của vị trí và vận tốc 
trong quá trình chuyển động.

\begin{figure}[H]
    \centering
    \includegraphics[width=0.8\textwidth]{gait_planner/cubic_pva.pdf}
    \caption{Các đường cong quỹ đạo bậc ba}
    \label{fig25}
\end{figure}

\subsection{Nội suy đa thức bậc năm}

Tuy nhiên, để cải thiện độ mượt của quỹ đạo và đảm bảo thêm điều kiện liên 
tục của gia tốc, phương pháp nội suy đa thức bậc năm được sử dụng. Dạng tổng quát:

\[
f(t) = a_5 t^5 + a_4 t^4 + a_3 t^3 + a_2 t^2 + a_1 t + a_0
\]

Các hệ số được xác định dựa trên sáu điều kiện biên gồm 
vị trí, vận tốc và gia tốc tại điểm đầu và điểm cuối:

\begin{align}
x(0) &= x_0, \quad x(T) = x_f \\
\dot{x}(0) &= 0, \quad \dot{x}(T) = 0 \\
\ddot{x}(0) &= 0, \quad \ddot{x}(T) = 0
\end{align}

Việc sử dụng đa thức bậc năm cho phép đảm bảo tính liên tục bậc hai ($C^2$) của vị trí, vận tốc và gia tốc, 
từ đó giúp chuyển động của robot trở nên ổn định và mượt mà hơn.

\begin{figure}[H]
    \centering
    \includegraphics[width=0.8\textwidth]{gait_planner/quintic_pva.pdf}
    \caption{Biên dạng vị trí, vận tốc và gia tốc của hàm bước mượt bậc năm}
    \label{fig26}
\end{figure}

Khi giải hệ sáu phương trình với sáu ẩn, nghiệm dạng đóng thu được dưới dạng tham số chuẩn hóa. Đặt $\tau = t/T \in [0,\,1]$, hàm bước mượt bậc năm (quintic smooth-step) được xác định là:

\begin{equation}
s(\tau) = 10\tau^3 - 15\tau^4 + 6\tau^5
\label{eq:quintic_step}
\end{equation}

Hàm $s(\tau)$ thỏa mãn tất cả sáu điều kiện biên:

\begin{equation}
s(0) = 0,\; s(1) = 1,\quad \dot{s}(0) = \dot{s}(1) = 0,\quad \ddot{s}(0) = \ddot{s}(1) = 0
\end{equation}

Nhờ đó, quỹ đạo giữa hai điểm bất kỳ có thể được viết gọn:

\begin{equation}
\mathbf{p}(\tau) = \mathbf{p}_0 + s(\tau) \cdot (\mathbf{p}_f - \mathbf{p}_0)
\end{equation}

trong đó $\mathbf{p}_0$ và $\mathbf{p}_f$ lần lượt là vị trí đầu và vị trí cuối.

\subsection{So sánh nội suy bậc ba và bậc năm}

Để làm rõ lý do lựa chọn phương pháp đa thức bậc năm thay vì bậc ba, ta xét biểu đồ so sánh biên dạng vị trí, vận tốc và gia tốc của cả hai phương pháp trên cùng một khoảng thời gian chuẩn hóa $\tau \in [0,\,1]$ (Hình~\ref{fig:cubic_vs_quintic}).

\begin{figure}[H]
    \centering
    \includegraphics[width=0.8\textwidth]{gait_planner/cubic_vs_quintic.pdf}
    \caption{So sánh biên dạng động học giữa nội suy bậc 3 và bậc 5.}
    \label{fig:cubic_vs_quintic}
\end{figure}

Từ Hình~\ref{fig:cubic_vs_quintic}, có thể thấy rõ sự khác biệt ở biên dạng gia tốc. Đối với đa thức bậc 3, gia tốc tại các thời điểm chuyển pha ($\tau=0$ và $\tau=1$) có giá trị khác không ($\pm 6$). Sự thay đổi đột ngột gia tốc từ 0 lên giá trị này gây ra lực giật (jerk), dẫn đến rung lắc cơ khí và lực va đập tại các khớp của robot. Ngược lại, đối với đa thức bậc 5, gia tốc tại $\tau=0$ và $\tau=1$ hoàn toàn bằng không, tạo sự chuyển tiếp mượt mà về mặt động lực học. Quỹ đạo đạt độ trơn bậc hai ($C^2$), hoàn toàn phù hợp với yêu cầu chuyển động ổn định của robot bốn chân trong thực tế.

Do đó, tác giả sử dụng hoàn toàn nội suy đa thức bậc năm cho việc tạo quỹ đạo bước chân nhằm tăng cường sự ổn định và kéo dài tuổi thọ của hệ thống truyền động.

\subsection{Triển khai quỹ đạo chân dạng chữ D}

Trong hệ thống thực tế, quỹ đạo của mỗi chân robot được chia thành hai pha: pha vung chân (swing phase) và pha chống chân (stance phase), tạo thành đường đi dạng chữ D trong không gian Descartes.

\textbf{\textit{Pha vung (Swing Phase)}} 

Trong pha vung, chân di chuyển từ vị trí bắt đầu $\mathbf{p}_A$ (phía sau) đến vị trí kết thúc $\mathbf{p}_D$ (phía trước) trong $N_{sw}$ bước. Tọa độ ngang ($x$, $y$) được nội suy bằng hàm bước mượt bậc năm:

\begin{align}
\tau &= \frac{i}{N_{sw} - 1}, \quad i = 0, 1, \ldots, N_{sw}-1 \\
s &= 10\tau^3 - 15\tau^4 + 6\tau^5 \\
x_i &= x_A + s \cdot (x_D - x_A) \\
y_i &= y_A + s \cdot (y_D - y_A)
\end{align}

Đối với tọa độ theo phương $z$, thay vì sử dụng đa thức bậc năm thuần túy, quỹ đạo nâng chân được thiết kế theo đường cong parabol trên biến $s$:

\begin{equation}
z_i = z_{stance} + \Delta h \cdot 4\,s\,(1 - s)
\label{eq:z_parabolic}
\end{equation}

trong đó $z_{stance}$ là chiều cao pha chống chân và $\Delta h$ là độ cao nâng chân. Biểu thức $4\,s\,(1-s)$ tạo ra đường cong hình chuông đối xứng đạt cực đại $\Delta h$ tại $s = 0.5$ (giữa pha vung chân), đồng thời trả về $z_{stance}$ tại $s = 0$ và $s = 1$ (đầu và cuối pha vung chân). Nhờ kế thừa tính liên tục $C^2$ từ hàm $s(\tau)$, quỹ đạo $z$ cũng mượt mà tại các điểm chuyển pha.

\textbf{\textit{Pha chống chân (Stance Phase)}} 

Trong pha chống chân, chân duy trì tiếp xúc mặt đất tại $z_{stance}$ và di chuyển ngược từ $\mathbf{p}_D$ về $\mathbf{p}_A$ trong $N_{st}$ bước, cũng sử dụng hàm bước mượt bậc năm:

\begin{align}
\tau &= \frac{i}{N_{st} - 1}, \quad i = 0, 1, \ldots, N_{st}-1 \\
s &= 10\tau^3 - 15\tau^4 + 6\tau^5 \\
x_i &= x_D + s \cdot (x_A - x_D) \\
y_i &= y_D + s \cdot (y_A - y_D) \\
z_i &= z_{stance}
\end{align}

Tổng số điểm quỹ đạo cho một chu kỳ dáng đi đầy đủ là $N_{total} = N_{sw} + N_{st}$.

\begin{figure}[H]
    \centering
    \includegraphics[width=0.8\textwidth]{gait_planner/gait_generator.png}
    \caption{Quỹ đạo chân 3D dạng chữ D (Pha vung Parabol và Pha chống đỡ đường thẳng)}
    \label{fig:gait_trajectory_3d}
\end{figure}

\subsection{Tham số quỹ đạo trong hệ thống thực tế}

Vị trí bắt đầu và kết thúc cho mỗi chân được tính theo:

\begin{align}
x_A &= x_{center} - \frac{L_{stride}}{2}, \quad x_D = x_{center} + \frac{L_{stride}}{2} \\
z_A &= z_D = z_{stance}
\end{align}

Hệ số chu kỳ nhiệm vụ (duty cycle) tổng quát được xác định là:

\begin{equation}
\beta = \frac{N_{st}}{N_{st} + N_{sw}}
\end{equation}

Dáng đi Trot yêu cầu giá trị $\beta > 0.5$, đảm bảo robot luôn có ít nhất hai chân tiếp xúc mặt đất tại mọi thời điểm để duy trì ổn định tĩnh trong suốt chu kỳ.

\subsection{Lập lịch pha dáng đi}

Dáng đi (gait) của robot bốn chân được xác định bởi mối quan hệ pha giữa các chân. Mỗi chân thực hiện cùng một quỹ đạo tuần hoàn, nhưng với độ lệch pha khác nhau. Pha dịch $\phi_i \in [0, 1)$ xác định vị trí tương đối của chân $i$ trong chu kỳ dáng đi tổng thể.

Trong triển khai thực tế, mỗi chân duy trì một chỉ số bước (waypoint index) $k_i$ trong mảng quỹ đạo tổng $N_{total} = N_{sw} + N_{st}$ điểm. Pha dịch được áp dụng bằng cách khởi tạo chỉ số bắt đầu của mỗi chân:

\begin{equation}
    k_{i,0} = \lfloor \phi_i \cdot N_{total} \rfloor
\end{equation}

Bảng~\ref{tab:gait_phases} trình bày các cấu hình pha dịch cho ba dáng đi phổ biến nhất.

\begin{table}[H]
\centering
\caption{Cấu hình pha dịch cho các dáng đi}
\label{tab:gait_phases}
\renewcommand{\arraystretch}{1.3}
\begin{tabular}{|l|c|c|c|c|c|}
\hline
\textbf{Dáng đi} & $\phi_{LF}$ & $\phi_{RF}$ & $\phi_{LB}$ & $\phi_{RB}$ & \textbf{Đặc điểm} \\ \hline
Trot & 0 & 0,5 & 0,5 & 0 & Hai chân chéo đồng pha \\ \hline
Walk & 0 & 0,5 & 0,75 & 0,25 & Tuần tự từng chân \\ \hline
Pace & 0 & 0,5 & 0 & 0,5 & Hai chân cùng bên đồng pha \\ \hline
\end{tabular}
\end{table}

\textbf{\textit{Dáng đi Trot}} 

($\phi = [0, 0.5, 0.5, 0]$) là cấu hình dáng đi phổ biến nhất. Hai chân chéo (LF--RB và RF--LB) hoạt động đồng pha, tạo ra hai nhóm chân xen kẽ. Tại mọi thời điểm, ít nhất hai chân thuộc cùng nhóm đang ở pha chống chân, duy trì ổn định tĩnh.

\textbf{\textit{Phân tích thời gian một chu kỳ}} 

Tổng thời gian của một chu kỳ dáng đi đầy đủ phụ thuộc vào số lượng khung hình (frame) của các pha và chu kỳ bộ định thời (timer period) $\Delta T$:

\begin{equation}
    T_{cycle} = (N_{sw} + N_{st}) \times \Delta T
\end{equation}

Tần số bước chân được tính theo chu kỳ:

\begin{equation}
    f_{step} = \frac{1}{T_{cycle}}
\end{equation}

Từ đó, tốc độ di chuyển tiến về phía trước theo lý thuyết là tích của chiều dài bước và tần số bước:

\begin{equation}
    v_{forward} = L_{stride} \times f_{step}
\end{equation}

Tốc độ này có thể được điều chỉnh trực tiếp thông qua việc thay đổi sải bước $L_{stride}$ hoặc thay đổi thời gian pha chống chân $N_{st}$.



\section{Thiết kế bộ điều khiển}

Để robot có thể duy trì thăng bằng và phản hồi chính xác với các nhiễu loạn ngoại lai (ví dụ: bề mặt địa hình gồ ghề), một hệ thống điều khiển vòng kín tự động dựa trên thuật toán PID được thiết kế. Cấu trúc tổng thể của hệ thống tự cân bằng này được trình bày chi tiết trong Hình~\ref{fig:applied_pid_simulink}.

Hệ thống hoạt động dựa trên cơ chế phản hồi trạng thái động học thông qua các khâu xử lý liên tiếp. Bắt đầu từ khâu đo lường và tạo sai số, cảm biến IMU (BNO055) liên tục đo lường góc nghiêng thực tế của thân robot (True Roll/Pitch). Các giá trị này được so sánh với góc tham chiếu lý thuyết (thường đặt $\phi_{ref}=0, \theta_{ref}=0$) thông qua các bộ cộng để tạo ra tín hiệu sai số $e_\phi$ và $e_\theta$. Tiếp đó, sai số được đưa vào bộ điều khiển PID (tích hợp khâu chống vọt lố Anti-windup và lọc thông thấp LPF), với đầu ra của khối này là các tín hiệu bù cân bằng (Compensation Signals). Tại khối chuyển mạch chế độ (Mode Switch), tùy thuộc vào trạng thái hiện tại của robot, hệ thống sẽ rẽ nhánh luồng tín hiệu bù. Đối với chế độ di chuyển (GAIT Mode), tín hiệu bù được gửi tới khối Quy hoạch quỹ đạo và Động học nghịch (Trajectory Planning \& Inverse Kinematics) để tính toán bù trừ trực tiếp vào tọa độ không gian làm việc của bàn chân. Ngược lại, đối với chế độ đứng yên (HOME Mode), tín hiệu bù được gửi tới khối Bù góc khớp trực tiếp (Direct Joint Angle Compensator) để tinh chỉnh không gian khớp. Cuối cùng, trong khâu thực thi, các góc khớp mục tiêu ($\theta_i$) được truyền xuống cơ cấu chấp hành (Bus Servo ST3215) để tác động lực lên đối tượng vật lý (Quadruped Robot Plant), từ đó tạo ra các chuyển động triệt tiêu nhiễu loạn và giúp thân robot lấy lại tư thế ngang bằng.

\begin{figure}[H]
    \centering
    \includegraphics[width=\textwidth]{img_matlab/fig-PID.png}
    \caption{Sơ đồ khối bộ điều khiển tự cân bằng của robot bốn chân.}
    \label{fig:applied_pid_simulink}
\end{figure}



\subsection{Kiến trúc điều khiển chi tiết}

Hệ thống điều khiển tổng thể của robot hoạt động với tần số vòng lặp 333~Hz ($\Delta T = 3$~ms), được thiết kế theo cấu trúc phân cấp chặt chẽ (Hình~\ref{fig:overall_architecture_locomotion}). Hệ thống tiếp nhận các lệnh cấp cao từ người dùng (User Commands) thông qua một máy trạng thái trung tâm (Command Manager) và điều phối luồng dữ liệu xuống các thành phần xử lý cốt lõi, bao gồm Bộ quy hoạch chuyển động (Locomotion Planner) và Bộ điều khiển cân bằng (Balance Controller).

\begin{figure}[H]
    \centering
    \includegraphics[width=\textwidth]{architecture/locomotion_balance_controller.png}
    \caption{Tổng quan kiến trúc điều khiển hệ thống robot bốn chân tích hợp bộ chuyển mạch chế độ tĩnh/động.}
    \label{fig:overall_architecture_locomotion}
\end{figure}

Đặc điểm nổi bật của kiến trúc này là khả năng phân chia luồng xử lý độc lập cho trạng thái di chuyển (Gait Mode) và trạng thái đứng yên (Home Mode), kết hợp với cơ chế phản hồi thăng bằng linh hoạt thông qua ba thành phần chính. Thành phần thứ nhất là Bộ quy hoạch chuyển động (Locomotion Planner), đóng vai trò là module chịu trách nhiệm tạo hình quỹ đạo chân khi robot di chuyển; module này bao gồm Bộ lập lịch pha (Gait Phase Scheduler) để đồng bộ thời gian các chân, kết hợp với Bộ tạo quỹ đạo (Trajectory Generator) để sinh ra tọa độ bước chân dạng đa thức Quintic và D-shape. Thành phần thứ hai là Khối bù tư thế không gian làm việc (Gait Attitude Compensator) nằm trong cụm Planner; khối này nhận tọa độ nguyên bản và cộng bù trực tiếp với tín hiệu chỉnh sửa $\Delta p$ từ bộ thăng bằng, sinh ra tọa độ chân hiệu chỉnh $p_{adj}$ trước khi đưa vào khâu giải động học nghịch (IK). Thành phần thứ ba là Bộ điều khiển cân bằng (Balance Controller) hoạt động hoàn toàn độc lập, sử dụng dữ liệu góc Roll/Pitch từ IMU làm phản hồi cho bộ PID để liên tục tính toán lượng bù chống lại ngoại lực gây lật thân.

Sự linh hoạt của hệ thống nằm ở công tắc chuyển mạch điều chỉnh cân bằng (\textit{Balance Output Switch}) và bộ chọn góc khớp cuối cùng (\textit{Mode-based Final Angle Selector}). Tùy thuộc vào chế độ hoạt động hiện tại, tín hiệu bù cân bằng sẽ được định tuyến theo hai phương thức hoàn toàn khác biệt. Trong chế độ tĩnh (HOME), hệ thống áp dụng cơ chế bù trong không gian khớp (Joint Space); khi robot đứng im, quỹ đạo làm việc bị khóa, tín hiệu bù từ PID được chuyển mạch thành độ dịch góc khớp ($\Delta\theta_3$). Độ dịch này sau đó được bộ Final Angle Selector cộng trực tiếp vào góc khớp gối danh định $\theta_{k\_home}$ để gửi ngay xuống phần cứng (UART Serial Driver). Việc bỏ qua hoàn toàn cụm tính toán IK giúp phản ứng cân bằng đạt độ trễ tối thiểu, khử nhiễu ngoại lực tức thời. Ngược lại, trong chế độ động (GAIT), hệ thống áp dụng cơ chế bù trong không gian làm việc (Task Space); khi robot đang bước đi, việc bảo toàn quỹ đạo chân 3D là yếu tố bắt buộc. Tín hiệu bù từ PID lúc này được chuyển đổi thành độ dời tọa độ ($\Delta p$) và đưa ngược vào bộ Gait Attitude Compensator. Bằng cách bù vào đầu vào của khối IK thay vì cộng trực tiếp vào góc khớp, robot có thể vừa giữ thăng bằng thân vừa bảo toàn chính xác hình dáng biên độ bước đi, tránh tình trạng trượt chân (slippage) tại các điểm đặt.

\subsection{Cơ chế áp dụng tín hiệu điều chỉnh tư thế}
Việc điều chỉnh tư thế được thực hiện theo thời gian thực thông qua quá trình đo đạc góc nghiêng từ cảm biến IMU. Đầu ra của các bộ điều khiển PID là các tín hiệu bù hiệu chỉnh $u_{roll}$ và $u_{pitch}$. Sự khác biệt cốt lõi giữa hai chế độ hoạt động nằm ở \textit{không gian} mà tín hiệu hiệu chỉnh được ánh xạ.

\subsubsection*{Áp dụng hiệu chỉnh trong không gian khớp ở chế độ tĩnh}
Khi robot ở trạng thái tĩnh (chế độ HOME), tín hiệu $u_{roll}, u_{pitch}$ từ bộ điều khiển PID được phân bổ thành giá trị độ lệch (offset) cho từng chân, với giới hạn bão hòa $|\Delta\theta_3| \leq \theta_{sat}$:
\begin{equation}
\begin{cases} 
\Delta\theta_{3,LF} = \mathrm{sat}\big(- (K_{home} u_{roll} - K_{home} u_{pitch}),\; \theta_{sat}\big) \\
\Delta\theta_{3,LB} = \mathrm{sat}\big(- (K_{home} u_{roll} + K_{home} u_{pitch}),\; \theta_{sat}\big) \\
\Delta\theta_{3,RF} = \mathrm{sat}\big(- (K_{home} u_{roll} + K_{home} u_{pitch}),\; \theta_{sat}\big) \\
\Delta\theta_{3,RB} = \mathrm{sat}\big(- (K_{home} u_{roll} - K_{home} u_{pitch}),\; \theta_{sat}\big) 
\end{cases}
\end{equation}
Lượng bù này được cộng trực tiếp làm giá trị hiệu chỉnh cho góc khớp gối ($\theta_3$). Do cấu trúc động học chuỗi hở, việc thay đổi trực tiếp góc khớp gối sẽ làm thay đổi vị trí bàn chân tương ứng. Động thái này tạo ra các phản lực từ mặt đất, giúp thân robot lấy lại trạng thái cân bằng tức thời mà không yêu cầu hệ thống phải tính toán qua khối động học nghịch (IK), qua đó giảm thiểu tối đa độ trễ điều khiển.

\subsubsection*{Áp dụng hiệu chỉnh trong không gian làm việc ở chế độ động}
Khi robot đang trong chu trình di chuyển (chế độ GAIT), biên dạng quỹ đạo bước chân bắt buộc phải được bảo toàn nghiêm ngặt để tránh hiện tượng trượt (slippage) và mất ổn định chu kỳ. Do đó, ma trận xoay nghịch đảo $R = R_y(u_{pitch}) R_x(u_{roll})$ được tính toán dựa trên lượng bù từ bộ điều khiển PID. Ma trận này sau đó được nhân với vectơ vị trí chân nguyên bản $\mathbf{p}_{raw}$ sinh ra từ bộ tạo quỹ đạo. Một cơ chế chọn lọc được áp dụng nhằm chỉ cập nhật duy nhất tọa độ cao độ trục Z của bàn chân, trong khi cố định hoàn toàn các giá trị tham chiếu trên trục X và Y:
\begin{equation}
\mathbf{p}_{adj} = R \cdot \mathbf{p}_{raw}, \quad \text{với} \quad x_{adj} = x_{raw}, \quad y_{adj} = y_{raw}
\end{equation}

Nguyên lý giữ nguyên sự biến thiên trên mặt phẳng XY đóng vai trò then chốt trong việc bảo toàn chiều dài sải bước (stride length) của robot. Tọa độ Đề-các hiệu chỉnh $[x_{adj}, y_{adj}, z_{adj}]$ sau đó mới được cấp vào bộ giải động học nghịch (IK) để suy ra góc quay mục tiêu cho các khớp. Quá trình biến đổi không gian này được minh họa trực quan trong Hình \ref{fig29}.

\begin{figure}[H]
    \centering
    \includegraphics[width=0.9\textwidth]{kinematics/attitude_adjustment.png}
    \caption{Minh họa quá trình điều chỉnh tư thế thân robot: từ $\mathbf{p}_{raw}$ (trước bù) sang $\mathbf{p}_{adj}$ (sau bù), chỉ thay đổi trục $Z$, giữ nguyên trục $X$ và $Y$.}
    \label{fig29}
\end{figure}

\subsection{Thuật toán điều khiển PID và các khâu bảo vệ}
Trong cả hai chế độ động và tĩnh, thuật toán cốt lõi đảm nhiệm triệt tiêu sai lệch tư thế là bộ điều khiển PID kinh điển, được phân tách tính toán độc lập cho trục Roll và trục Pitch. Đầu vào của khối là sai số $e(t) = 0 - \theta_{meas}(t)$, trong đó $\theta_{meas}$ là giá trị góc nghiêng đo được từ IMU sau khi đã đi qua bộ lọc trung bình trượt hàm mũ (EMA) nhằm khử nhiễu cơ học tần số cao.

Luật điều khiển PID được triển khai dưới dạng phương trình sai phân rời rạc như sau:
\begin{equation}
    u_{PID}(k) = K_P e(k) + K_I \sum_{i=0}^{k} e(i) \Delta T + K_D \frac{e(k) - e(k-1)}{\Delta T}
\end{equation}

Nhằm đáp ứng yêu cầu an toàn tuyệt đối cho phần cứng robot khi vận hành thực tế, tín hiệu đầu ra của bộ PID buộc phải đi qua các lớp bảo vệ thiết yếu trước khi được phát lệnh tới cơ cấu chấp hành. Lớp bảo vệ đầu tiên là khâu triệt vùng chết (Deadband), chủ động bỏ qua các vi sai quá nhỏ dao động quanh vị trí cân bằng ($0$ rad); kỹ thuật này triệt tiêu hiện tượng chattering (động cơ rung lắc liên tục), qua đó giảm hao mòn cơ khí và tỏa nhiệt không cần thiết. Lớp bảo vệ thứ hai bao gồm bão hòa tín hiệu (Saturation) và lọc thông thấp (LPF); lệnh điều khiển $u_{PID}$ được giới hạn biên độ trong khoảng an toàn $[-u_{max}, u_{max}]$ để ngăn chặn các xung bù vượt quá khả năng chịu tải của cơ hệ, sau đó tín hiệu được làm mượt qua bộ lọc thông thấp (Low-Pass Filter) để các khớp chuyển động êm ái, hạn chế tối đa các lực giật (jerk) cục bộ. Lớp bảo vệ cuối cùng là cơ chế chống tích lũy (Anti-windup); khi cơ cấu chấp hành hoặc tín hiệu đầu ra chạm ngưỡng bão hòa vật lý, cơ chế anti-windup sẽ tự động tạm dừng hoặc suy giảm thành phần tích phân $I$. Điều này chặn đứng sự tích lũy sai số ảo dư thừa, giúp hệ thống phục hồi nhanh chóng và không bị vọt lố (overshoot) khi các nhiễu loạn ngoại lực kết thúc.

Việc phân bổ linh hoạt chiến lược bù theo không gian làm việc tương ứng với trạng thái động học, kết hợp cùng các cơ chế bảo vệ tín hiệu nhiều lớp, là cơ sở khoa học then chốt giúp hệ thống điều khiển vận hành ổn định và bám sát lý thuyết động lực học.
